\lab{4}{Лемматизация и стемминг}
\task
Добавить поиск без учёта словоформ; сравнить поведение с точным поиском,
рассмотреть ухудшившиеся запросы и возможный бонус за точное совпадение.
Описать выбранную нормализацию и оценку качества после её включения.

\method
Реализованы три режима: \code{exact}, \code{stem} и \code{lemma}.
Первый оставляет токены после понижения регистра. Второй использует
самостоятельную реализацию английского стеммера Портера. Третий применяет
словарь WordNet 3.0 и собственные правила получения кандидатов.
Одинаковая функция нормализации применяется при индексации и поиске.

Портер последовательно преобразует суффиксы, учитывая наличие гласной,
меру чередования гласных и согласных, удвоенные согласные и окончания CVC.
Результат не обязан быть словом языка: \code{connections} даёт
\code{connect}. Слишком короткие слова и токены с символами вне
\code{a--z} не меняются. Поэтому \code{pthread_create},
\code{std::vector}, числа и русские токены не повреждаются.

Словарный метод сначала проверяет исключения и наличие исходной формы,
затем строит кандидатов по окончаниям для существительных, глаголов,
прилагательных и наречий. Замена принимается только при наличии кандидата
в словаре с нужным признаком части речи. Так обрабатываются нерегулярные
формы \code{children -> child}, \code{mice -> mouse},
\code{went -> go}. Контекст предложения и снятие омонимии не используются.
При неоднозначности важен фиксированный порядок правил.

WordNet используется как набор данных, а не как готовая поисковая
библиотека. Скрипт \code{build_lexicon.py} упаковывает словарные записи
и исключения; CMake встраивает массив в исполняемый файл.
Собственный бинарный поиск работает по упакованному словарю.
Идентификатор словаря записывается в координатный индекс, чтобы исключить
его чтение с несовместимой нормализацией.

\journal
Сначала реализована защита идентификаторов и каскад Портера, затем
добавлены словарь, исключения и морфологические правила WordNet.
Для сравнения сохранён режим без морфологии. В булевом индексаторе
предусмотрен дополнительный exact-индекс: сначала выводятся документы
с точной формой, затем остальные совпадения по основе. В штатной
конфигурации дополнительный индекс отключён. Лабораторные 7--8 используют
стемминг, 10--13 и 17 - WordNet.

\testing
Проверяются классические примеры шагов Портера, нерегулярные формы,
слова без преобразований, пустая строка и сохранение идентификаторов.
Особый случай - увеличение длины при \code{mice -> mouse}:
при достаточном буфере преобразование проходит, при недостаточном
должна возникать ошибка. Для качества нужно сравнить одинаковые запросы
и корпус в трёх режимах, разметить объединённый пул и оценить метрики
из лабораторной №6. Совпадение основы само по себе не доказывает
релевантность найденного документа.

\results
\begin{table}[H]\centering
\begin{tabular}{lll}\toprule
Вход & Портер & WordNet\\\midrule
\code{connections} & \code{connect} & \code{connection}\\
\code{allocations} & \code{alloc} & \code{allocation}\\
\code{children} & \code{children} & \code{child}\\
\code{running} & \code{run} & \code{run}\\
\code{pthread_create} & \code{pthread_create} & \code{pthread_create}\\\bottomrule
\end{tabular}
\caption{Ожидаемые преобразования, задаваемые правилами и тестами.}
\end{table}
В полном булевом индексе со стеммингом 19\,787\,654 терма, в
координатном индексе с WordNet - 20\,430\,767. Разница составляет
643\,113 термов (около 3,15\% от словаря WordNet), но это сравнение
двух способов нормализации. Оно не измеряет сокращение относительно
отдельно построенного exact-индекса.

В выдаче по полному корпусу запросы \code{memory allocation} и
\code{memory allocations} дают одинаковые первые результаты TF-IDF
и по 282\,185 кандидатов: формы объединяются. Для запроса
\code{epolll} исправление опечатки не выполняется и результатов нет.
Оценка качества системы с WordNet приведена в лабораторной №6.
Изолированное влияние морфологии требует отдельного сравнения
exact и lemma при одинаковом отборе и ранжировании кандидатов.
Сравнение строгого поиска и TF-IDF такой эксперимент не заменяет.

\conclusions
Нормализация объединяет словоформы, но может объединять разные смыслы
и оставлять часть родственных слов раздельными. Словарь полезен для
нерегулярных форм, однако не решает проблему контекста и опечаток.
Для технических терминов сохранение точной формы и аккуратное расширение
запроса могут быть полезнее агрессивного сведения к основе.
Изолированное сравнение режимов следует проводить на общем пуле
с оценками релевантности для каждого варианта.
